% Section IX: Contrasts and limits - what is and is not topological
% Purpose: Clarify distinction between "topological effects in QFT" and "TQFT proper"
% Unify lecture material (topological sectors) with formal program (TQFT)
%
% Key contrasts to make explicit:
% - Chern-Simons IS topological
% - Maxwell-Chern-Simons is NOT
% - Yang-Mills with θ-term is NOT
% - Instantons are topological sectors but YM is not TQFT
% - Monopoles have topological charge but theory is not TQFT
%
% Status: EXTRACT AND ORGANIZE from existing comparative discussions
% Sources:
% - tex_docs/anomalies_boundaries_topological_response.tex (Maxwell-CS contrast, lines ~500-600)
% - writings/chern_simons_review.tex (comparative discussions throughout)
% - Section II content (once written) will provide setup for these contrasts
%
% Action needed: Write new section synthesizing the contrasts
% Use Maxwell-CS as primary worked example (most detailed in sources)
% Add brief discussions of YM + θ, instantons, monopoles

\section{What is and is not a Topological Quantum Field Theory}

% TODO: Write synthesis section clarifying boundaries

\subsection{Metric Independence as the Defining Feature}

% TODO: Define what makes a QFT "topological"
% - Correlation functions invariant under smooth metric deformations: δZ/δg_μν = 0
% - No local propagating degrees of freedom (or they decouple)
% - Observables are topological invariants
%
% Contrast: ordinary QFT
% - Metric appears essentially (Hodge star, kinetic terms)
% - Local propagating modes (particles, wavelengths, energies)
% - Even when there are topological SECTORS, the dynamics are not topological

\subsection{Chern--Simons Theory: A Genuine TQFT}

% TODO: Review CS theory against TQFT criteria
% - Action has no metric: S = (k/4π)∫tr(A∧dA + ⅔A∧A∧A)
% - F = 0 equation of motion: no propagating gluons
% - Wilson loops are topological (knot/link invariants)
% - Partition function Z(M) is topological invariant
%
% This is why it's TQFT: topology controls the physics

\subsection{Maxwell--Chern--Simons Theory: Not Topological}

% TODO: Extract and expand from sources
% Source: anomalies_boundaries_topological_response.tex, lines ~500-600
% Also: chern_simons_review.tex Maxwell-CS discussion
%
% Action: S = ∫[(1/4)F∧*F + (k/4π)ε^{μνρ}A_μ∂_νA_ρ]
% Hodge star *F requires metric
% Propagating photon (even though massive from CS term)
% Useful effective theory for FQHE edges, BUT not a TQFT
%
% Lesson: Adding CS term to ordinary gauge theory doesn't make it topological

% FIGURE FLAG [F-mcs-propagator]: Comparison of CS vs MCS propagator structure

\subsection{\texorpdfstring{Yang--Mills with a $\theta$-Term}{Yang--Mills with a theta-Term}}

% TODO: S_YM + θ∫F∧F̃
% θ-term is topological (total derivative), labels sectors
% BUT Yang-Mills itself still has metric, gluon dynamics
% Topological STRUCTURE within non-topological theory

\subsection{Instantons and Monopoles: Topological Features vs. TQFT}

% TODO: Connect back to Section II (topological structures in ordinary QFT)
%
% Instantons: configuration space has π₃(G) sectors, ∫F∧F̃ is topological
% But YM action still has local dynamics
%
% Monopoles: magnetic charge from π₂(G/H), Dirac quantization
% But gauge-Higgs theory still has propagating fields
%
% These are topological OBJECTS in ordinary QFT, not TQFTs

\subsection{Summary: The Distinction}

% TODO: Table or clear summary
% Topological structure inside QFT:
%   - sectors, charges, defects classified by topology
%   - theory still has local dynamics and metric dependence
%
% Topological QFT:
%   - topology itself controls observables
%   - metric-independent (or decoupled)
%   - no propagating local degrees of freedom (or only topological sector)
%
% Slogan: "Ordinary QFT + topological sectors ≠ TQFT"
